\begin{frame}{그림으로 읽기: 하드 마진과의 세 가지 차이}
  \begin{columns}[T]
    \begin{column}{0.52\textwidth}
      \centering
      \includegraphics[width=\linewidth]{ch05_2_soft_margin.png}
    \end{column}
    \begin{column}{0.44\textwidth}
      \small
      \begin{enumerate}
        \item \kb{마진 경계 안에 데이터가 들어와 있다} — 5.1절에선
          $f=\pm1$ 사이에 아무 점도 없었지만, 여기서는 이상치가
          마진 안(또는 경계 너머)에 있다. 허용된 것이지 결함이
          아니다.
        \item \kb{서포트 벡터의 위치가 바뀌었다} — 이상치 두 개가
          서포트 벡터에 포함(코드 출력: 인덱스 80, 81 모두
          $f=+1.000$, $\xi=0$). 이상치를 ``받아안으려다''
          마진 경계 위에 정확히 올려놓은 것.
        \item \kb{마진이 5.1절보다 좁아졌다} — $2/\|w\|$가
          $\approx8.5$보다 작다. 이상치를 받으려면 $\|w\|$가 커져
          마진이 좁아진다 — $C$가 조절하는 트레이드오프의 단면.
      \end{enumerate}
      \vskip0.4em
      {\footnotesize 이상치는 ``완벽한 이상치''(다른 클래스 근처)가
      아니라 \textbf{경계 쪽 클러스터 사이에 놓인} 점. $C=1$에선
      어느 영역(② 또는 ③)에 빠지는지 예측해라.}
    \end{column}
  \end{columns}
\end{frame}
